\section{Global radial-curvature certificate for the Xi kernel}
\label{sec:globalcurvature}

The monotonicity of the potential curvature for the
classical Riemann theta kernel is not new: Newman proved
in 1991 that $V'$ is convex on the positive half-line,
using a GHS-type argument from statistical mechanics
\cite{NewmanGHS}. Our aim here is different and narrower:
to give a reproducible quantitative verification by
outward-rounded rational intervals and an explicit
infinite-range estimate, while isolating the exact
bounds subsequently used in the Laguerre comparison.
In particular, this is not a claim of first discovery
of the underlying monotonic-curvature property.
The numerical and analytic certificate still requires
independent specialist review.

\begin{theorem}[Global monotone theta curvature]
\label{thm:globalcurvature}
For the classical Jacobi-theta kernel
\begin{align*}
\Phi(u)&=\sum_{m\ge1}\pi m^2
  \bigl(2\pi m^2e^{4u}-3\bigr)\\
&\qquad\cdot e^{5u-\pi m^2e^{4u}}.
\end{align*}
let $V(u)=-\log\Phi(u)$. Then
\begin{equation}\label{eq:globalthird}
\boxed{V'''(u)>0\quad\text{for every }u>0.}
\end{equation}
Consequently $V''$ is strictly increasing in
$|u|$. In the variable $t=u^2$, the function
\[
t^{3/2}\frac{\mathcal Ts(t)}{s(t)^2}
\quad (s(t)=\Phi(\sqrt t))
\]
is strictly increasing on $(0,\infty)$.
\end{theorem}

\paragraph{Exact reduction.}
Put $x=\pi e^{4u}$ and
\begin{align}\label{eq:Aj}
A_j(x)&=\sum_{m\ge1}(-1)^j m^{2+2j}
 \bigl(2m^2x-3-2j\bigr)\nonumber\\
&\qquad\cdot e^{-(m^2-1)x},
\quad 0\le j\le4.
\end{align}
Then
$\Phi(u)=\pi e^{5u-x}A_0(x)$ and
$A_j'(x)=A_j(x)+A_{j+1}(x)$ for $j\le3$.
All $A_j$ are absolutely convergent and
smooth for $x\ge\pi$. Since $A_0(x)>0$,
define the cubic numerator
\begin{align}
H(x)={}&-x^3
\bigl(A_3A_0^2-3A_1A_2A_0+2A_1^3\bigr)
\nonumber\\
&-3x^2\bigl(A_2A_0^2-A_1^2A_0\bigr)
-xA_1A_0^2.
\label{eq:H}
\end{align}
Three applications of $4x\partial_x$ to
$\log A_0(x)-x+5u+\log\pi$ give the
\emph{exact identity}
\begin{equation}\label{eq:thirdH}
V'''(u)=\frac{64H(x)}{A_0(x)^3}.
\end{equation}
The Jacobi functional equation makes $\Phi$
even about $u=0$, so $V'''(0)=0$ and
$H(\pi)=0$.

\paragraph{Finite interval portion.}
The standalone standard-library script
\texttt{certificates/xi\_global\_curvature\_exact.py}
works with common denominator $10^{65}$,
rounding every multiplication, reciprocal,
and integer division outward. A rigorous
Machin-series enclosure bounds $\pi$.
For $0\le z<1/4$, an 80-term Taylor
expansion encloses $e^z$, since
\[
0<e^z-\sum_{k=0}^{80}\frac{z^k}{k!}
 \le \frac{2}{4^{81}81!}<10^{-65}.
\]
Exponentials $e^{-y}$ are then enclosed by
\[
e^{-y}=\bigl(e^{y/4096}\bigr)^{-4096}.
\]
For every $x\in[\pi,6]$, truncating
\eqref{eq:Aj} at $m=8$ causes an error,
for each $0\le j\le4$, smaller than
\[
2\cdot23\cdot9^{12}\cdot2^{-240}
<10^{-58}.
\]
Indeed the $m$th absolute contribution
for $m\ge9$ is bounded by
$23m^{12}2^{-3(m^2-1)}$, and its
successive ratio is less than $1/2$.
The script explicitly enlarges all five
intervals by $10^{-58}$.

Using $A_j'=A_j+A_{j+1}$ and
forward automatic differentiation of
\eqref{eq:H}, $100$ rational interval cells
covering $[\pi,7/2]$ establish
\begin{equation}
\boxed{H'(x)>46\quad(\pi\le x\le7/2).}
\label{eq:Hprimebound}
\end{equation}
With $H(\pi)=0$, this yields $H(x)>0$
for $\pi<x\le7/2$. Independently,
$400$ rational interval cells on $[7/2,6]$
establish
\begin{equation}
\boxed{H(x)>26\quad(7/2\le x\le6).}
\label{eq:Hcompactbound}
\end{equation}
These are outward enclosures on whole
intervals, not tests only at mesh nodes.

\paragraph{Analytic infinite tail.}
For $x\ge6$, let
\[
a_j=(-1)^j(2x-3-2j),\qquad E_j=A_j-a_j.
\]
For $0\le j\le3$ and $m\ge2$,
the $m$th summand of $E_j$ has magnitude
at most $3xm^{10}e^{-(m^2-1)x}$.
Successive terms have ratio bounded by
$(3/2)^{10}e^{-30}<1/2$.
Consequently
\begin{equation}\label{eq:Etail}
|E_j(x)|\le6144x e^{-3x}=:\epsilon(x)x,
\qquad 0\le j\le3.
\end{equation}
The first-summand numerator, a purely
algebraic simplification, is
\begin{equation}\label{eq:H0}
H_0(x)=x(8x^3-36x^2+42x-45)
>\frac74x^4
\end{equation}
for $x\ge6$. The latter strict inequality
follows, for $t=x-6\ge0$, from
\begin{align*}
H_0(x)-\frac74x^4
={}&\frac{25}{4}t^4+114t^3+744t^2\\
&+1971t+1566>0.
\end{align*}
Along every straight-line interpolation
between $(a_0,\ldots,a_3)$ and
$(A_0,\ldots,A_3)$, we have
$|A_j|\le3x$ because $\epsilon(x)<1$.
Bounding the four first partial derivatives
of the cubic \eqref{eq:H} by the
multivariate mean-value inequality yields
\begin{equation}\label{eq:Herror}
|H-H_0|\le \epsilon(x)x^4
(162x^2+162x+27).
\end{equation}
The function
$e^{-3x}(162x^2+162x+27)$ is
strictly decreasing for $x\ge6$.
Finally the elementary rational Taylor bound
$e^3>20$ gives $e^{18}>20^6=64{,}000{,}000$;
hence for every $x\ge6$,
\[
\frac{|H-H_0|}{x^4}
<\frac{6144\cdot6831}{64{,}000{,}000}
<\frac74.
\]
Equations \eqref{eq:H0}--\eqref{eq:Herror}
give $H(x)>0$ on $[6,\infty)$.
Combining the three ranges with
\eqref{eq:thirdH} proves
Theorem~\ref{thm:globalcurvature}.

\begin{proposition}[Unconditional common Laguerre band]
\label{prop:commonlaguerre}
Let $F(z)=\int_\mathbb{R}\Phi(u)e^{izu}\,du$,
and let $A_n,B_n$ be the positive
correlation moments of
Section~\ref{sec:transformattack}.
For every integer $n\ge1$,
\[
\frac{B_n}{A_n}\le\frac{B_1}{A_1}.
\]
In particular,
\[
\boxed{L_n[F](x)>0\quad
\text{for all }n\ge1,\quad
|x|<\sqrt{2A_1/B_1}.}
\]
\end{proposition}
\begin{proof}
Theorem~\ref{thm:globalcurvature}
establishes the radial monotonicity
hypothesis of
Proposition~\ref{prop:radial}.
The reverse-likelihood-ratio argument
therefore proves the moment-ratio
inequality. Applying
$L_n[F](x)\ge A_n-\tfrac12x^2B_n$
finishes the proof.
\end{proof}

\paragraph{Certification and limitations.}
The two finite inequalities above are
established by exact integer operations and
explicit analytic theta/exponential tails;
the entire proof is independently auditable
from the standalone Python 3 script.
The algebraic identities and the rational
infinite-tail margin are included in the
Lean module
\texttt{XiKernel/GlobalCurvatureAlgebra.lean}.
The \emph{full} theta-series positivity
statement and the interval verifier have
not yet been formally imported and proved
within Lean's kernel.
Neither the global radial-curvature theorem
nor its common finite Fourier-positivity
band implies the Riemann hypothesis.
Newman's prior theorem \cite{NewmanGHS} establishes
the principal qualitative convexity assertion.
Any possible contribution of this section is the
alternative constructive certificate, its strict
positive rational margins on the stated domains,
and the way those estimates are applied here;
their novelty has not been independently established.
Independent specialist review remains outstanding.
